% Entiers quadratiques : équation de Pell-Fermat et polygones réguliers du réseau
% Compile avec pdflatex, xelatex ou lualatex.
\documentclass[11pt,a4paper]{article}

\usepackage{iftex}
\ifPDFTeX
  \usepackage[utf8]{inputenc}
  \usepackage[T1]{fontenc}
  \IfFileExists{lmodern.sty}{\usepackage{lmodern}}{}
  \usepackage[french]{babel}
\else
  \usepackage{fontspec}
  \usepackage{polyglossia}
  \setdefaultlanguage{french}
\fi
\usepackage[margin=2.6cm]{geometry}
\usepackage{amsmath,amssymb,amsthm,mathtools}
\usepackage{enumitem}
\usepackage{microtype}
\usepackage[hidelinks]{hyperref}

% ---------- Environnements ----------
\theoremstyle{plain}
\newtheorem{theoreme}{Théorème}[section]
\newtheorem{proposition}[theoreme]{Proposition}
\newtheorem{lemme}[theoreme]{Lemme}
\newtheorem{corollaire}[theoreme]{Corollaire}
\theoremstyle{definition}
\newtheorem{definition}[theoreme]{Définition}
\newtheorem{exemple}[theoreme]{Exemple}
\newtheorem{exercice}{Exercice}
\theoremstyle{remark}
\newtheorem{remarque}[theoreme]{Remarque}

% ---------- Notations ----------
\newcommand{\N}{\mathbb{N}}
\newcommand{\Z}{\mathbb{Z}}
\newcommand{\Q}{\mathbb{Q}}
\newcommand{\R}{\mathbb{R}}
\newcommand{\C}{\mathbb{C}}
\newcommand{\pgcd}{\wedge}
\DeclareMathOperator{\Nm}{N}
\DeclareMathOperator{\tr}{tr}

\title{Entiers quadratiques : unités de $\Z[\sqrt2]$\\ et polygones réguliers du réseau\\[2pt]
\large Exercices commentés pour l'agrégation}
\author{Michel Hua}
\date{Octobre 2026}

\begin{document}
\maketitle

\begin{abstract}
On reprend deux exercices de terminale. Le premier montre qu'aucun triangle équilatéral n'a ses sommets sur un quadrillage, par l'irrationalité de $\sqrt3$. Le second établit que, pour tout $q\in\N$, $(1+\sqrt2)^q=\sqrt{n+1}+\sqrt n$ pour un entier $n$. Derrière le second se trouve la structure du groupe des unités de $\Z[\sqrt d]$ : on montre que les unités de $\Z[\sqrt2]$ sont les $\pm(1+\sqrt2)^k$, ce qui résout les équations $x^2-2y^2=\pm1$, et que $\Z[\sqrt2]$ est euclidien. Derrière le premier se trouve un énoncé plus fort : un polygone régulier à $n$ sommets dans $\Z^2$ n'existe que pour $n=4$, et dans $\Z^d$ que pour $n\in\{3,4,6\}$. La preuve repose sur un argument de trace : une rotation qui préserve un réseau a une trace entière.
\end{abstract}

\paragraph{Leçons concernées.} 126 (Exemples d'équations en arithmétique), 122 (Anneaux principaux), 191 (Exemples d'utilisation des techniques d'algèbre en géométrie), 102 (Groupe des nombres complexes de module $1$, racines de l'unité).

\paragraph{Développements possibles.} Théorème \ref{thm:unites} et corollaire \ref{cor:z2} (unités de $\Z[\sqrt2]$ et équation $x^2-2y^2=\pm1$) ; proposition \ref{prop:euclide} et exercice \ref{exo:premiers} ; théorème \ref{thm:polygones} avec l'exercice \ref{exo:zd}.

\paragraph{Notations.} Dans tout le texte, $d\geq2$ est un entier qui n'est pas un carré parfait. On note $\Z[\sqrt d]=\{a+b\sqrt d : a,b\in\Z\}$ et $\Q(\sqrt d)=\{r+s\sqrt d : r,s\in\Q\}$, vus comme sous-anneaux de $\R$.

\tableofcontents

% ==========================================================
\section{Unités de \texorpdfstring{$\Z[\sqrt d]$}{Z[sqrt d]} et équation de Pell-Fermat}
% ==========================================================

\subsection{Norme et conjugaison}

\begin{lemme}\label{lem:irr}
Le réel $\sqrt d$ est irrationnel. Par conséquent, si $r,s\in\Q$ vérifient $r+s\sqrt d=0$, alors $r=s=0$, et l'écriture d'un élément de $\Q(\sqrt d)$ sous la forme $r+s\sqrt d$ est unique.
\end{lemme}

\begin{proof}
Supposons $\sqrt d=p/q$ avec $p,q\geq1$ et $p\pgcd q=1$. Alors $p^2=dq^2$. Si un nombre premier $\ell$ divisait $q$, il diviserait $p^2$, donc $p$, ce qui contredit $p\pgcd q=1$. Ainsi $q=1$ et $d=p^2$ est un carré, ce qui est exclu. Si $r+s\sqrt d=0$ avec $s\neq0$, alors $\sqrt d=-r/s\in\Q$, ce qui est absurde. Donc $s=0$, puis $r=0$. L'unicité de l'écriture s'en déduit par différence.
\end{proof}

\begin{definition}
La \emph{conjugaison} est l'application $\sigma:\Q(\sqrt d)\to\Q(\sqrt d)$, $r+s\sqrt d\mapsto r-s\sqrt d$, bien définie par le lemme \ref{lem:irr}. La \emph{norme} est $\Nm(x)=x\,\sigma(x)$, c'est-à-dire $\Nm(r+s\sqrt d)=r^2-ds^2$.
\end{definition}

\begin{proposition}\label{prop:norme}
\begin{enumerate}[nosep]
\item $\sigma$ est un automorphisme involutif de l'anneau $\Q(\sqrt d)$, qui stabilise $\Z[\sqrt d]$.
\item La norme est multiplicative, prend des valeurs entières sur $\Z[\sqrt d]$, et $\Nm(x)=0$ si et seulement si $x=0$.
\item Un élément $x\in\Z[\sqrt d]$ est inversible dans $\Z[\sqrt d]$ si et seulement si $\Nm(x)=\pm1$, et alors $x^{-1}=\Nm(x)\,\sigma(x)$.
\end{enumerate}
\end{proposition}

\begin{proof}
1. L'additivité est claire. Pour $x=r+s\sqrt d$ et $x'=r'+s'\sqrt d$, on a $xx'=(rr'+dss')+(rs'+r's)\sqrt d$, et
\[
  \sigma(x)\sigma(x')=(r-s\sqrt d)(r'-s'\sqrt d)=(rr'+dss')-(rs'+r's)\sqrt d=\sigma(xx').
\]
Enfin $\sigma(1)=1$, $\sigma\circ\sigma=\mathrm{id}$, et $\sigma$ envoie $\Z[\sqrt d]$ dans lui-même.

2. $\Nm(xx')=xx'\sigma(x)\sigma(x')=\Nm(x)\Nm(x')$. Si $x\in\Z[\sqrt d]$, $\Nm(x)=a^2-db^2\in\Z$. Si $\Nm(r+s\sqrt d)=0$ avec $s\neq0$, alors $d=(r/s)^2$ et $\sqrt d=|r/s|\in\Q$, ce qui contredit le lemme \ref{lem:irr} ; donc $s=0$, puis $r=0$.

3. Si $xy=1$ avec $x,y\in\Z[\sqrt d]$, alors $\Nm(x)\Nm(y)=1$ dans $\Z$, donc $\Nm(x)=\pm1$. Réciproquement, si $\Nm(x)=\pm1$, alors $x\cdot\Nm(x)\sigma(x)=\Nm(x)^2=1$ et $\Nm(x)\sigma(x)\in\Z[\sqrt d]$.
\end{proof}

Ainsi, résoudre l'équation de Pell-Fermat $x^2-dy^2=\pm1$ dans $\Z^2$, c'est déterminer le groupe $\Z[\sqrt d]^\times$.

\subsection{Structure du groupe des unités}

\begin{lemme}\label{lem:positif}
Si $u=a+b\sqrt d\in\Z[\sqrt d]^\times$ vérifie $u>1$, alors $a\geq1$ et $b\geq1$. En particulier $u\geq1+\sqrt d$.
\end{lemme}

\begin{proof}
Par la proposition \ref{prop:norme}, $\sigma(u)=\Nm(u)/u=\pm1/u\in\,]-1,1[$. Donc $2a=u+\sigma(u)>1-1=0$ et $2b\sqrt d=u-\sigma(u)>1-1=0$. Comme $a$ et $b$ sont entiers, $a\geq1$ et $b\geq1$.
\end{proof}

\begin{theoreme}\label{thm:unites}
Si $\Z[\sqrt d]^\times\neq\{\pm1\}$, alors $\Z[\sqrt d]^\times$ possède un plus petit élément $\varepsilon_d$ strictement supérieur à $1$, et
\[
  \Z[\sqrt d]^\times=\{\pm\varepsilon_d^{\,k} : k\in\Z\}\simeq\{\pm1\}\times\Z .
\]
On dit que $\varepsilon_d$ est l'\emph{unité fondamentale} de $\Z[\sqrt d]$.
\end{theoreme}

\begin{proof}
\emph{Existence de $\varepsilon_d$.} Soit $u$ une unité différente de $\pm1$. Alors $|u|\neq1$ puisque $u$ est réel, et l'un des quatre réels $u,-u,u^{-1},-u^{-1}$, qui sont des unités, est $>1$ ; notons-le $M$. Par le lemme \ref{lem:positif}, une unité $a+b\sqrt d\in\,]1,M]$ vérifie $1\leq a\leq M$ et $1\leq b\leq M/\sqrt d$ ; il n'y en a donc qu'un nombre fini, et au moins une, $M$. On note $\varepsilon_d$ la plus petite.

\emph{Descente.} Soit $u$ une unité. Quitte à remplacer $u$ par $-u$, on peut supposer $u>0$. Comme $\varepsilon_d>1$, il existe un unique $k\in\Z$ tel que $\varepsilon_d^{\,k}\leq u<\varepsilon_d^{\,k+1}$, à savoir $k=\lfloor\ln u/\ln\varepsilon_d\rfloor$. Alors $v=u\,\varepsilon_d^{-k}$ est une unité vérifiant $1\leq v<\varepsilon_d$. Par minimalité de $\varepsilon_d$, on ne peut avoir $v>1$, donc $v=1$ et $u=\varepsilon_d^{\,k}$.

Enfin $(\pm1,k)\mapsto\pm\varepsilon_d^{\,k}$ est un morphisme de groupes surjectif, injectif car $\pm\varepsilon_d^{\,k}=1$ impose $\varepsilon_d^{\,k}=1$, donc $k=0$, et le signe $+$.
\end{proof}

\begin{corollaire}\label{cor:z2}
\begin{enumerate}[nosep]
\item $\Z[\sqrt2]^\times=\{\pm(1+\sqrt2)^k : k\in\Z\}$.
\item Pour $q\in\N$, écrivons $(1+\sqrt2)^q=a_q+b_q\sqrt2$. Les solutions dans $\N^2$ de $x^2-2y^2=1$ sont les $(a_q,b_q)$ avec $q$ pair, et celles de $x^2-2y^2=-1$ sont les $(a_q,b_q)$ avec $q$ impair.
\end{enumerate}
\end{corollaire}

\begin{proof}
1. $\Nm(1+\sqrt2)=1-2=-1$, donc $1+\sqrt2$ est une unité $>1$ et le théorème s'applique. Par le lemme \ref{lem:positif}, toute unité $>1$ est $\geq1+\sqrt2$, donc $\varepsilon_2=1+\sqrt2$.

2. Si $(x,y)\in\N^2$ vérifie $x^2-2y^2=\pm1$, alors $u=x+y\sqrt2$ est une unité, non nulle, donc $(x,y)\neq(0,0)$ et $u\geq1$. Par le point 1, $u=(1+\sqrt2)^q$ avec $q\in\Z$, et $q\geq0$ car $(1+\sqrt2)^q<1$ pour $q<0$. Par unicité de l'écriture, $(x,y)=(a_q,b_q)$. Enfin $a_q^2-2b_q^2=\Nm\big((1+\sqrt2)^q\big)=(-1)^q$. Réciproquement chaque $(a_q,b_q)$ est dans $\N^2$, par la récurrence $a_{q+1}=a_q+2b_q$, $b_{q+1}=a_q+b_q$, $(a_0,b_0)=(1,0)$.
\end{proof}

\begin{exemple}
Les premières valeurs sont, pour $q=0,\dots,7$ :
\[
\begin{array}{c|cccccccc}
q&0&1&2&3&4&5&6&7\\\hline
a_q&1&1&3&7&17&41&99&239\\
b_q&0&1&2&5&12&29&70&169
\end{array}
\]
Les deux suites vérifient $x_{q+2}=2x_{q+1}+x_q$, puisque $1+\sqrt2$ est racine de $X^2-2X-1$.
\end{exemple}

\subsection{L'exercice d'origine}

L'exercice demandait de montrer que, pour tout $q\in\N$, il existe $n\in\N$ tel que $(1+\sqrt2)^q=\sqrt{n+1}+\sqrt n$. La feuille d'origine montrait par deux récurrences parallèles que $(1-\sqrt2)^q=a_q-b_q\sqrt2$. C'est exactement l'énoncé $\sigma\big((1+\sqrt2)^q\big)=\sigma(1+\sqrt2)^q$, qui découle du fait que $\sigma$ est un morphisme d'anneaux. On obtient ainsi un énoncé général.

\begin{proposition}\label{prop:consecutifs}
Soit $(a,b)\in\N^2$ tel que $a^2-db^2=\eta\in\{-1,1\}$. Pour $q\in\N$, on écrit $(a+b\sqrt d)^q=A_q+B_q\sqrt d$. Alors $A_q,B_q\in\N$ et $A_q^2-dB_q^2=\eta^q$. En particulier les entiers $A_q^2$ et $dB_q^2$ sont consécutifs, et
\[
  (a+b\sqrt d)^q=\sqrt{n+1}+\sqrt n,\qquad n=\min(A_q^2,\,dB_q^2),
\]
cet entier $n$ étant unique.
\end{proposition}

\begin{proof}
Par récurrence, $A_0=1$, $B_0=0$, et $A_{q+1}=aA_q+dbB_q$, $B_{q+1}=bA_q+aB_q$ sont des entiers positifs. Par multiplicativité de la norme, $A_q^2-dB_q^2=\Nm\big((a+b\sqrt d)^q\big)=\Nm(a+b\sqrt d)^q=\eta^q$. Comme $A_q\geq0$ et $B_q\geq0$, on a $A_q=\sqrt{A_q^2}$ et $B_q\sqrt d=\sqrt{dB_q^2}$, deux racines d'entiers consécutifs. L'unicité vient de ce que $t\mapsto\sqrt{t+1}+\sqrt t$ est strictement croissante sur $\R_+$.
\end{proof}

\begin{exemple}
Pour $d=2$ et $a=b=1$, on trouve $n=2b_q^2$ si $q$ est pair et $n=a_q^2$ si $q$ est impair. Ainsi $(1+\sqrt2)^5=41+29\sqrt2=\sqrt{1681}+\sqrt{1682}$ et $(1+\sqrt2)^6=99+70\sqrt2=\sqrt{9801}+\sqrt{9800}$.
\end{exemple}

\begin{remarque}
La preuve n'utilise pas que $d$ n'est pas un carré si l'on remplace $\sigma$ par le développement du binôme : $(a-b\sqrt d)^q=A_q-B_q\sqrt d$, car seuls les termes d'exposant impair en $\sqrt d$ changent de signe.
\end{remarque}

\subsection{Compléments}

\begin{proposition}\label{prop:euclide}
L'anneau $\Z[\sqrt2]$ est euclidien pour le stathme $x\mapsto|\Nm(x)|$. En particulier il est principal, donc factoriel.
\end{proposition}

\begin{proof}
L'anneau $\Z[\sqrt2]$ est intègre, comme sous-anneau de $\R$. Soit $x,y\in\Z[\sqrt2]$ avec $y\neq0$. Dans $\Q(\sqrt2)$, on écrit $x/y=x\sigma(y)/\Nm(y)=r+s\sqrt2$ avec $r,s\in\Q$. On choisit $m,n\in\Z$ avec $|r-m|\leq\frac12$ et $|s-n|\leq\frac12$, et on pose $q=m+n\sqrt2$ et $\rho=x-qy\in\Z[\sqrt2]$. Avec $\alpha=r-m$ et $\beta=s-n$, on a $\rho=y(\alpha+\beta\sqrt2)$, donc
\[
  |\Nm(\rho)|=|\Nm(y)|\cdot|\alpha^2-2\beta^2|\leq|\Nm(y)|\max(\alpha^2,2\beta^2)\leq\tfrac12|\Nm(y)|<|\Nm(y)|,
\]
la première inégalité venant de ce que $\alpha^2$ et $2\beta^2$ sont positifs.
\end{proof}

\begin{remarque}
Le même argument montre que $\Z[\sqrt3]$ est euclidien, avec $|\alpha^2-3\beta^2|\leq\frac34$. Il ne s'étend pas à tout $d$ : $\Z[\sqrt{10}]$ n'est pas principal. On a en effet $2\cdot3=(4+\sqrt{10})(4-\sqrt{10})$, et $2$ est irréductible sans être premier. Il est irréductible car aucun élément n'est de norme $\pm2$, puisque $x^2\equiv\pm2\pmod5$ n'a pas de solution. Il n'est pas premier car il ne divise ni $4+\sqrt{10}$ ni $4-\sqrt{10}$.
\end{remarque}

\begin{remarque}[Fractions continues]
On a $\sqrt2=[1;2,2,2,\dots]$, et les réduites $1,\frac32,\frac75,\frac{17}{12},\frac{41}{29},\dots$ sont exactement les $a_q/b_q$. On le retrouve ainsi : si $a^2-2b^2=\pm1$ avec $b\geq1$, alors $a\geq b$ et
\[
  \Big|\frac ab-\sqrt2\Big|=\frac{1}{b\,(a+b\sqrt2)}<\frac{1}{2b^2},
\]
donc $a/b$ est une réduite de $\sqrt2$ par le critère de Legendre. En général, l'unité fondamentale de $\Z[\sqrt d]$ se lit sur la période du développement en fraction continue de $\sqrt d$.
\end{remarque}

\begin{remarque}[Dirichlet]
Le théorème \ref{thm:unites} ne dit pas que $\Z[\sqrt d]^\times\neq\{\pm1\}$. C'est vrai pour tout $d$ (Lagrange) : l'équation $x^2-dy^2=1$ a toujours une solution non triviale, ce qui se démontre par le théorème d'approximation de Dirichlet. C'est le cas particulier des corps quadratiques réels du théorème des unités de Dirichlet : pour un corps de nombres $K$ ayant $r_1$ plongements réels et $r_2$ paires de plongements complexes, $\mathcal O_K^\times\simeq\mu(K)\times\Z^{r_1+r_2-1}$. Notons que $\Z[\sqrt d]$ n'est pas toujours l'anneau des entiers de $\Q(\sqrt d)$ : pour $d=5$, celui-ci est $\Z\big[\frac{1+\sqrt5}2\big]$, d'unité fondamentale le nombre d'or.
\end{remarque}

% ==========================================================
\section{Polygones réguliers à sommets entiers}
% ==========================================================

\subsection{Le triangle équilatéral}

\begin{proposition}\label{prop:triangle}
Il n'existe pas de triangle équilatéral, non réduit à un point, dont les sommets sont dans $\Z^2$.
\end{proposition}

\begin{proof}[Preuve de la feuille d'origine]
On identifie $\R^2$ à $\C$. Soit $A,B,C$ trois points de $\Z^2=\Z[i]$ formant un triangle équilatéral. Quitte à échanger $B$ et $C$, il est direct, ce qui s'écrit $c-a=\omega(b-a)$ avec $\omega=e^{i\pi/3}=\frac12+i\frac{\sqrt3}2$. Posons $b-a=u+iv$ et $c-a=s+it$ avec $u,v,s,t\in\Z$. En séparant parties réelle et imaginaire,
\[
  2s-u=-v\sqrt3,\qquad 2t-v=u\sqrt3 .
\]
Par le lemme \ref{lem:irr} appliqué à $d=3$, on obtient $u=v=0$, donc $B=A$, puis $C=A$.
\end{proof}

\begin{proof}[Preuve par l'aire]
Soit $A,B,C\in\Z^2$ formant un triangle équilatéral de côté $\ell>0$. D'une part, son aire vaut $\frac12|\det(B-A,C-A)|\in\frac12\Z$. D'autre part, elle vaut $\frac{\sqrt3}{4}\ell^2$ avec $\ell^2=\|B-A\|^2\in\Z$ non nul. Donc $\sqrt3=2|\det(B-A,C-A)|/\ell^2\in\Q$, ce qui est absurde.
\end{proof}

\subsection{Le cas général}

\begin{definition}
Soit $n\geq3$ et $P$ un plan affine euclidien. Un \emph{polygone régulier à $n$ sommets} de $P$ est une famille $(P_k)_{k\in\Z/n\Z}$ de la forme $P_k=O+R^k(v)$, où $O\in P$, $v$ est un vecteur non nul de la direction $V$ de $P$ et $R$ est la rotation vectorielle de $V$ d'angle $2\pi/n$ pour l'une des deux orientations.
\end{definition}

\begin{theoreme}\label{thm:polygones}
Soit $n\geq3$. Il existe un polygone régulier à $n$ sommets dont les sommets sont dans $\Z^2$ si et seulement si $n=4$.
\end{theoreme}

\begin{proof}
Le carré $(0,0),(1,0),(1,1),(0,1)$ convient pour $n=4$. Réciproquement, soit $(P_k)$ un tel polygone, $P_k=O+R^k(v)$, et soit $\theta=2\pi/n\in\,]0,2\pi/3]$.

\emph{Les vecteurs côtés.} Posons $w_k=P_{k+1}-P_k=R^k(Rv-v)$. Ce sont des vecteurs de $\Z^2$ comme différences de points de $\Z^2$, et $R(w_k)=w_{k+1}$ pour tout $k\in\Z/n\Z$. Le vecteur $w_0=Rv-v$ est non nul, car $R$ n'a pas d'autre point fixe que $0$ puisque $\theta\notin2\pi\Z$. Les vecteurs $w_0$ et $w_1=Rw_0$ sont linéairement indépendants, car $\theta\notin\pi\Z$.

\emph{Le réseau.} Soit $L=\sum_{k}\Z w_k$, le sous-groupe de $\Z^2$ engendré par les $w_k$. Comme $R$ permute les $w_k$, on a $R(L)=L$. Sous-groupe de $\Z^2$, $L$ est un $\Z$-module libre de rang fini. Des vecteurs de $\Z^2$ linéairement indépendants sur $\Z$ le sont sur $\Q$ (on chasse les dénominateurs), donc sur $\R$ (le rang d'une matrice rationnelle ne dépend pas du corps). Le rang de $L$ est donc au plus $2$, et il vaut $2$ puisque $L$ contient $w_0$ et $w_1$. Une $\Z$-base $(e_1,e_2)$ de $L$ est ainsi une base de $\R^2$.

\emph{La trace.} Comme $R(e_1),R(e_2)\in L$, la matrice de $R$ dans la base $(e_1,e_2)$ est à coefficients entiers, donc $\tr R=2\cos\theta\in\Z$. Or $\theta\in\,]0,2\pi/3]$ donne $2\cos\theta\in[-1,2[$, donc $2\cos\theta\in\{-1,0,1\}$ et $\theta\in\{2\pi/3,\pi/2,\pi/3\}$, c'est-à-dire $n\in\{3,4,6\}$.

\emph{Conclusion.} Le cas $n=3$ est exclu par la proposition \ref{prop:triangle}. Si $n=6$, les points $P_0,P_2,P_4$ s'écrivent $O+(R^2)^j(v)$, $j=0,1,2$, avec $R^2$ de rotation d'angle $2\pi/3$ : ils forment un triangle équilatéral à sommets dans $\Z^2$, ce qui est exclu. Il reste $n=4$.
\end{proof}

\begin{remarque}
La preuve vaut aussi pour les polygones étoilés, où $R$ est la rotation d'angle $2\pi m/n$ avec $m\pgcd n=1$ : la condition $2\cos(2\pi m/n)\in\{-1,0,1\}$ donne encore $n\in\{3,4,6\}$. Elle reste valable mot pour mot dans $\Z^d$ (exercice \ref{exo:zd}) : seule l'exclusion de $n=3$ et $n=6$ utilise la dimension $2$.
\end{remarque}

\begin{remarque}[Racines de l'unité]
En termes complexes, si $P_k=z_k\in\Z[i]$, alors $e^{2i\pi/n}=(z_2-z_1)/(z_1-z_0)\in\Q(i)$. Le théorème \ref{thm:polygones} équivaut donc à : les racines de l'unité de $\Q(i)$ sont $\pm1,\pm i$. On peut aussi le voir ainsi : une racine primitive $n$-ième de l'unité est de degré $\varphi(n)$ sur $\Q$ (irréductibilité des polynômes cyclotomiques). Si elle est dans $\Q(i)$, on a $\varphi(n)\leq2$, donc $n\in\{1,2,3,4,6\}$, et $e^{2i\pi/3}=\frac{-1+i\sqrt3}2\notin\Q(i)$.
\end{remarque}

\begin{remarque}[Entiers d'Eisenstein]
Dans le réseau hexagonal $\Z[j]$, $j=e^{2i\pi/3}$, il y a des triangles équilatéraux et des hexagones réguliers, mais pas de carré, puisque $i\notin\Q(j)=\Q(i\sqrt3)$. Le plan $x+y+z=0$ coupe $\Z^3$ selon une copie de ce réseau, ce qui explique l'exercice \ref{exo:zd}.
\end{remarque}

% ==========================================================
\section{Exercices}
% ==========================================================

\begin{exercice}\label{exo:puissances}
Calculer $(1+\sqrt2)^q$ pour $1\leq q\leq5$, et l'écrire sous la forme $\sqrt{n+1}+\sqrt n$.
\end{exercice}

\begin{exercice}\label{exo:z3}
Montrer que l'unité fondamentale de $\Z[\sqrt3]$ est $2+\sqrt3$, que toutes les unités de $\Z[\sqrt3]$ sont de norme $1$, et résoudre $x^2-3y^2=1$ dans $\N^2$.
\end{exercice}

\begin{exercice}\label{exo:triangulaires}
Déterminer les nombres triangulaires $\frac{n(n+1)}2$, $n\geq1$, qui sont des carrés parfaits.
\end{exercice}

\begin{exercice}\label{exo:premiers}
Soit $p$ un nombre premier tel que $2$ soit un carré modulo $p$. Montrer qu'il existe $(x,y)\in\Z^2$ tel que $p=x^2-2y^2$. Traiter $p=7$ et $p=17$.
\end{exercice}

\begin{exercice}\label{exo:zd}
Soit $d\geq2$. Montrer qu'un polygone régulier à $n$ sommets dont les sommets sont dans $\Z^d$ n'existe que pour $n\in\{3,4,6\}$, et que les trois cas sont réalisés dans $\Z^3$.
\end{exercice}

% ==========================================================
\section{Corrigés}
% ==========================================================

\paragraph{Exercice \ref{exo:puissances}.} Par la récurrence $a_{q+1}=a_q+2b_q$, $b_{q+1}=a_q+b_q$ :
\[
\begin{aligned}
1+\sqrt2&=\sqrt2+\sqrt1, & 3+2\sqrt2&=\sqrt9+\sqrt8, & 7+5\sqrt2&=\sqrt{50}+\sqrt{49},\\
17+12\sqrt2&=\sqrt{289}+\sqrt{288}, & 41+29\sqrt2&=\sqrt{1682}+\sqrt{1681}.
\end{aligned}
\]
Les valeurs de $n$ sont donc $1,8,49,288,1681$, conformément à la proposition \ref{prop:consecutifs}.

\paragraph{Exercice \ref{exo:z3}.} On a $\Nm(2+\sqrt3)=4-3=1$, donc $2+\sqrt3$ est une unité $>1$. Soit $u=a+b\sqrt3$ une unité $>1$ ; par le lemme \ref{lem:positif}, $a,b\geq1$. Si $b=1$, alors $a^2=3\pm1$, donc $a=2$. Si $b\geq2$, alors $u\geq1+2\sqrt3>2+\sqrt3$, puisque $\sqrt3>1$. Donc $\varepsilon_3=2+\sqrt3$, et $\Z[\sqrt3]^\times=\{\pm(2+\sqrt3)^k\}$ par le théorème \ref{thm:unites}. Toutes ces unités sont de norme $\Nm(2+\sqrt3)^k=1$. On le voit aussi directement : $x^2-3y^2=-1$ imposerait $x^2\equiv2\pmod3$, ce qui est impossible. Comme dans le corollaire \ref{cor:z2}, les solutions de $x^2-3y^2=1$ dans $\N^2$ sont les coefficients de $(2+\sqrt3)^k$, $k\geq0$ : $(1,0)$, $(2,1)$, $(7,4)$, $(26,15)$, $(97,56)$, \dots

\paragraph{Exercice \ref{exo:triangulaires}.} On a $\frac{n(n+1)}2=m^2\iff(2n+1)^2-2(2m)^2=1$. Posons $x=2n+1$ et $y=2m$. Réciproquement, si $x^2-2y^2=1$, alors $x$ est impair, donc $x^2\equiv1\pmod8$, donc $8\mid2y^2$ et $y$ est pair. Par le corollaire \ref{cor:z2}, les solutions avec $n\geq1$ correspondent à $x+y\sqrt2=(1+\sqrt2)^{2k}=(3+2\sqrt2)^k$, $k\geq1$ :
\[
  (x,y)=(3,2),\ (17,12),\ (99,70),\ (577,408),\dots\qquad
  (n,m)=(1,1),\ (8,6),\ (49,35),\ (288,204),\dots
\]
Par exemple $\frac{49\cdot50}2=1225=35^2$.

\paragraph{Exercice \ref{exo:premiers}.} Soit $a\in\Z$ avec $p\mid a^2-2=(a-\sqrt2)(a+\sqrt2)$. L'élément $p$ n'est pas inversible dans $\Z[\sqrt2]$, car $\Nm(p)=p^2$. Il ne divise ni $a-\sqrt2$ ni $a+\sqrt2$ dans $\Z[\sqrt2]$ : on aurait $a\pm\sqrt2=p(u+v\sqrt2)$, donc $pv=\pm1$. Ainsi $p$ n'est pas premier dans $\Z[\sqrt2]$. Cet anneau étant principal (proposition \ref{prop:euclide}), $p$ n'est donc pas irréductible : $p=\pi\rho$ avec $\pi,\rho$ non inversibles. Alors $\Nm(\pi)\Nm(\rho)=p^2$ avec $|\Nm(\pi)|,|\Nm(\rho)|\neq1$, donc $\Nm(\pi)=\pm p$. En écrivant $\pi=x+y\sqrt2$, on a $x^2-2y^2=\pm p$. Si le signe est $-$, on remplace $\pi$ par $\pi(1+\sqrt2)$, de norme $-\Nm(\pi)=p$. Pour $p=7$, $3^2\equiv2$ et $7=3^2-2\cdot1^2$. Pour $p=17$, $6^2\equiv2$ et $17=5^2-2\cdot2^2$.

Inversement, si $p=x^2-2y^2$ est impair, alors $x$ est impair, et $p\equiv1-2y^2\equiv1$ ou $7\pmod 8$. Avec la seconde loi complémentaire de réciprocité quadratique, $p$ impair est de la forme $x^2-2y^2$ si et seulement si $p\equiv\pm1\pmod8$.

\paragraph{Exercice \ref{exo:zd}.} Soit $(P_k)$ un polygone régulier à $n$ sommets dans $\Z^d$, contenu dans un plan affine de direction $V$, et soit $R$ la rotation de $V$ correspondante. On reprend la preuve du théorème \ref{thm:polygones}. Les vecteurs $w_k=P_{k+1}-P_k$ sont dans $\Z^d\cap V$, et $R$ les permute. Le groupe $L$ qu'ils engendrent est un sous-groupe de $\Z^d$, donc libre de rang fini. Une $\Z$-base de $L$ est formée de vecteurs entiers $\Z$-libres, donc $\R$-libres, contenus dans $V$ : le rang est au plus $2$, et il vaut $2$ puisque $w_0$ et $w_1$ sont indépendants. On conclut comme dans le théorème que $2\cos(2\pi/n)\in\Z$, donc $n\in\{3,4,6\}$.

Dans $\Z^3$, le carré $(0,0,0),(1,0,0),(1,1,0),(0,1,0)$ convient pour $n=4$, et le triangle $(1,0,0),(0,1,0),(0,0,1)$, de côté $\sqrt2$, pour $n=3$. Pour $n=6$, on prend les points
\[
  (1,-1,0),\ (1,0,-1),\ (0,1,-1),\ (-1,1,0),\ (-1,0,1),\ (0,-1,1),
\]
c'est-à-dire $\rho^k(1,-1,0)$, $0\leq k\leq5$, avec $\rho(x,y,z)=(-y,-z,-x)$. L'application $\rho$ est une isométrie linéaire de $\R^3$ qui stabilise le plan $V=\{x+y+z=0\}$. On a $\rho^2(x,y,z)=(z,x,y)$ et $\rho^3=-\mathrm{id}$, donc la restriction de $\rho$ à $V$ est d'ordre exactement $6$, puisque ni $\rho^2$ ni $\rho^3$ n'y induisent l'identité. Une isométrie d'ordre $6$ d'un plan euclidien est une rotation d'angle $\pm\pi/3$, puisque les symétries orthogonales sont d'ordre $2$. Les six points forment donc un hexagone régulier de centre $0$ et de rayon $\sqrt2$. Dans ce cas, la preuve par l'aire de la proposition \ref{prop:triangle} ne s'applique plus : l'aire du triangle $(1,0,0),(0,1,0),(0,0,1)$ vaut $\frac{\sqrt3}2\notin\frac12\Z$.

\begin{thebibliography}{9}
\bibitem{samuel} P. Samuel, \emph{Théorie algébrique des nombres}, Hermann.
\bibitem{duverney} D. Duverney, \emph{Théorie des nombres}, Dunod.
\bibitem{hw} G. H. Hardy, E. M. Wright, \emph{An Introduction to the Theory of Numbers}, Oxford University Press.
\bibitem{ir} K. Ireland, M. Rosen, \emph{A Classical Introduction to Modern Number Theory}, Springer.
\bibitem{gourdon} X. Gourdon, \emph{Les maths en tête. Algèbre et probabilités}, Ellipses.
\end{thebibliography}

\paragraph{Source.} Feuille d'exercices de terminale S (2013).

\end{document}
